% PRISM: Photonic Similarity Engine for KV Cache Block Selection in Long-Context LLM Inference
% Format: Physical Review Letters (REVTeX 4.2)
% Authors: Hyoseok Park, Yeonsang Park
% Affiliation: Department of Physics, Chungnam National University

\documentclass[pra,twocolumn,superscriptaddress,amsmath,amssymb,longbibliography,floatfix]{revtex4-2}

% ── Packages ──────────────────────────────────────────────
\usepackage{graphicx}
\usepackage[colorlinks=true,linkcolor=blue,citecolor=blue,urlcolor=blue]{hyperref}
\usepackage{cleveref}
% Fix revtex4-2 + cleveref compatibility: revtex4-2 internally uses
% \section* (unnumbered sections), causing cleveref to write an empty
% label type to the .aux file ("cref reference format for label type
% '' undefined").  Define cleveref formats for the empty-string type
% so the lookup \csname cref@@format\endcsname succeeds.
\makeatletter
\expandafter\gdef\csname cref@@format\endcsname#1#2#3{#2Sec.~#1#3}
\expandafter\gdef\csname Cref@@format\endcsname#1#2#3{#2Section~#1#3}
\expandafter\gdef\csname crefrange@@format\endcsname#1#2#3#4#5#6{%
  Secs.~#1#2 to~#4#5#6}
\expandafter\gdef\csname Crefrange@@format\endcsname#1#2#3#4#5#6{%
  Sections~#1#2 to~#4#5#6}
\makeatother
\crefformat{section}{#2Sec.~#1#3}
\Crefformat{section}{#2Section~#1#3}
\crefname{section}{Sec.}{Secs.}
\Crefname{section}{Section}{Sections}
\usepackage{booktabs}
\usepackage{siunitx}
\usepackage{multirow}
\usepackage{tabularx}
\usepackage{xcolor}
\usepackage{enumitem}
\usepackage{placeins}   % \FloatBarrier to prevent figures drifting to end
\usepackage{bibunits}
\defaultbibliographystyle{apsrev4-2}

% ── Custom commands ───────────────────────────────────────
\newcommand{\vect}[1]{\mathbf{#1}}
\newcommand{\mat}[1]{\mathbf{#1}}
\newcommand{\TODO}[1]{\textcolor{red}{\textbf{[TODO: #1]}}}
\newcommand{\prism}{\textsc{Prism}}

% ══════════════════════════════════════════════════════════
\begin{document}

% ── Title ─────────────────────────────────────────────────
\title{PRISM: Photonic Similarity Engine for KV Cache Block Selection\\in Long-Context LLM Inference}

\author{Hyoseok Park}
\affiliation{Department of Physics, Chungnam National University, Daejeon 34134, Republic of Korea}

\author{Yeonsang Park}
\affiliation{Department of Physics, Chungnam National University, Daejeon 34134, Republic of Korea}

\date{\today}

% ── Abstract ──────────────────────────────────────────────
\begin{abstract}
Long-context LLM inference is bottlenecked not by compute
but by the memory bandwidth required to scan the KV cache
at every decode step---a cost that grows linearly with context
length.
The semiconductor industry increasingly acknowledges this shift:
NVIDIA's Vera Rubin architecture dedicates an entire DPU (ICMS)
to KV cache management with flash-backed storage and
hardware-assisted prefetch---an architectural bet confirming that
memory, not arithmetic, is the first-class
constraint.

Recent photonic accelerators have demonstrated impressive
throughput for dense attention
computation.
However, these approaches inherit the same $O(n)$ memory scaling
as electronic attention when applied to long contexts.
We observe that the real leverage point is the coarse
block-selection step: a memory-bound similarity search that
determines which KV blocks to fetch.
We identify, for the first time, that this task is
\emph{structurally matched} to the photonic broadcast-and-weight
paradigm---the query fans out to all candidates via passive
splitting, signatures are quasi-static (matching electro-optic
MRR programming), and only rank order matters (relaxing precision
to 4--6 bits).
Crucially, the photonic advantage \emph{grows with context
length}: as $N$ increases, the electronic scan cost rises
linearly while the photonic evaluation remains $O(1)$.

We instantiate this insight in \prism{}, a thin-film lithium
niobate (TFLN) similarity engine.
Hardware-impaired needle-in-a-haystack evaluation on Qwen2.5-7B
confirms 100\% accuracy from 4K through 64K tokens at $k{=}32$,
with $32\times$ traffic reduction.
\prism{} achieves a four-order-of-magnitude energy advantage
over GPU baselines at practical context lengths ($n \geq 4$K).
\end{abstract}

\maketitle

% ══════════════════════════════════════════════════════════
% MAIN TEXT (bibunit 1)
% ══════════════════════════════════════════════════════════
\begin{bibunit}

% ── Section files ─────────────────────────────────────────
\input{sections/introduction}           % 1. Introduction
\input{sections/background}             % 2. Background
\input{sections/architecture}           % 3. Photonic Retrieval Architecture
\input{sections/hardware_analysis}      % 4. Photonic Hardware Analysis
\input{sections/system_evaluation}      % 5. System-Level Evaluation
\input{sections/scaling_analysis}       % 6. Photonic Scaling Analysis
\input{sections/discussion}             % 7. Discussion
\input{sections/conclusion}             % 8. Conclusion

% ── References (main text) ───────────────────────────────
\putbib[references]

\end{bibunit}

\clearpage  % Force supplementary to start on a new page

% ══════════════════════════════════════════════════════════
% SUPPLEMENTARY (bibunit 2)
% ══════════════════════════════════════════════════════════
\begin{bibunit}

\input{sections/supplementary}

% ── References (supplementary) ───────────────────────────
\putbib[references]

\end{bibunit}

\end{document}
